\documentclass[3p, number, 11pt, sort&compress]{elsarticle}

\usepackage{amsmath}
\usepackage{tgtermes}
\usepackage[utf8]{inputenc}
\usepackage[T1]{fontenc}
\usepackage[labelfont=bf, justification=justified, singlelinecheck=false]{caption}
\usepackage{multicol,multirow,array}
\usepackage{lineno}
\linenumbers
\emergencystretch 3em
\usepackage{float}


\title{}

\author[1]{Emilio Bianchi\corref{cor1}}
\ead{ebianchi@unrn.edu.ar}

\author[2]{Tomás Guozden}
 
\author[3]{Juan Rivera}

 \affiliation[1]{organization={CONICET - Universidad Nacional de R\'io Negro},
 	country={Argentina}} 
 \affiliation[2]{organization={Universidad Nacional de R\'io Negro},
        city={Buenos Aires},
        country={Argentina}}
 \affiliation[3]{organization={CONICET - Universidad de Buenos Aires. Centro de Investigaciones del Mar y la Atm\'osfera (CIMA)},
        country={Argentina}}

        
\cortext[cor1]{Corresponding author. E-mail address: ebianchi@unrn.edu.ar}

\begin{document}

\begin{abstract}

\end{abstract}

\begin{keyword}
 \sep  \sep  \sep 
\end{keyword}

\maketitle

\section{Introduction}

Power systems are becoming increasingly sensitive to weather variations \cite{hoste2009matching}. Both electricity demand and supply are affected by atmospheric conditions and its variations through time and space \cite{dubus2018does, staffell2018increasing}. Electricity demand is mainly affected by temperature \cite{moral2005modelling, de2015seasonal, cassarino2018impact}, as well as by other weather variables such as wind speed, humidity, solar radiation and cloud cover \cite{psiloglou2009factors, apadula2012relationships}. Power generation from weather-dependent renewable sources such as wind and solar is determined by the availability of wind kinetic energy and solar radiation plus temperature, respectively \cite{de2019scoping, bloomfield2021sub}. These  variables fluctuate over different timescales, from micro-meteorological through synoptic, seasonal and up to decadal variations \cite{troccoli2012long,dubus2018does, staffell2018increasing, lledo2019seasonal}; and affect many stakeholders from the electricity market: facility owners, operation and maintenance teams, energy traders and system administrators \cite{bloomfield2021sub}. Anticipating climatic conditions from multiple timescales is crucial for decision making processes and energy management at local and national scales \cite{de2015seasonal, soret2019sub}.

The implementation of short-range weather forecasts is a common practice in the energy sector to support intra-daily to weekly operations\cite{son2017short, drew2017importance, browell2018improved, stanger2019optimising, de2019scoping}. Despite the fact that weather and climate variations influence decisions made at longer timescales \cite{kirschen2018fundamentals}, the implementation of sub-seasonal to seasonal forecasts has been limited so far \cite{bloomfield2021sub}. This could be due to the overall lower skill of these forecasts (when compared to synoptic forecasts), plus higher complexity and associated uncertainty \cite{bruno2016barriers, lledo2019seasonal, bloomfield2021sub}. Recently, both sub-seasonal to seasonal ranges have gained attention due to an overall improvement  of climate predictions \cite{clark2017skilful, lledo2019seasonal,merryfield2020current, meehl2021initialized}, which are now being used to anticipate power demand \cite{clark2017skilful}, wind and solar power generation \cite{lledo2019seasonal, soret2019sub, bett2022simplified}.

One of the main drivers of sub-seasonal climate variability worldwide is the Madden Julian Oscilation (MJO) \cite{zhang2005madden, zhang2013madden, kim2018prediction}. The MJO is an organized pattern of eastward-propagating tropical convection with a life cycle of about 40–50 days \cite{zhang2005madden, kim2018prediction} that affect weather variables and atmospheric circulation not only in the tropics, but also at mid-latitudes \cite{donald2006near} through the excitation of stationary Rossby waves \cite{seo2012global, zheng2018impacts, chen2021diversity}. A number of studies have reported relationships between this oscillation and weather variables and circulation patterns over different regions worldwide \cite{donald2006near, lin2009influence, matsueda2015global, zheng2018impacts}. In fact, is though to be an important contributor predictability in the extratropics due to recent advances in forecast techniques and skill \cite{waliser2003potential, matsueda2015global, kim2018prediction}. The predictability horizon of the MJO (up to one month) is longer than the temporal reach of synoptic forecasts, and this can provide valuable information for the management of climate-sensitive systems \cite{kim2014predictability, neena2014predictability, oliver2016predictability}. Nevertheless, studies referring to the MJO as potential predictor of renewable energy resources are so far scarce  \cite{lledo2020predicting}.

I propose to explore the statistical relationships between the different phases of the MJO and its amplitude and renewable energy resources (wind speed and solar radiation) over different regions of interest worldwide, using measured (when available), and reanalysis-derived data. 


\section{Data and methodology} \label{sec:dataandmothods}



\section{Results} \label{sec:results}



\section{Conclusions} \label{sec:conclusions}




\end{document}

